\section{Related Work}

Selective prediction lets a model decline to answer. The classical formulation
weighs the error rate on accepted inputs against the rejection rate
\citep{chow1970optimum}, and the statistical foundations of the trade-off are
well studied \citep{elyaniv2010foundations,bartlett2008classification,cortes2016learning}.
Modern instantiations attach a confidence score and threshold it
\citep{geifman2017selective,hendrycks2017baseline,gangrade2021selective}, and
selective answering is now standard for question answering and vision-language
models under domain shift \citep{kamath2020selective,whitehead2022reliable}.
This literature decides whether to answer a query. Our decision is different in
two ways: the system chooses a subset of auxiliary inputs rather than answering
directly, and the quantity that can be harmed is the prediction of a fixed
downstream model, not the label of the classifier that produces the score.

Distribution-free calibration provides the machinery we use. Conformal
prediction turns a score into a set or interval with finite-sample coverage
\citep{vovk2005algorithmic,lei2018distribution,angelopoulos2023gentle}, and its
risk-controlling variants bound the expectation of a monotone loss instead of
interval coverage \citep{bates2021distribution,angelopoulos2023conformal}.
Quantile and false-positive variants extend the same idea to other objectives
\citep{snell2023quantile,fisch2022conformal}. We adopt the risk-control view
because the operational requirement in context selection is a budget on
harmful selections, not coverage of a utility interval. The paper contributes
the instantiation of that view to context routing, the comparison against
interval certification in the regime where it collapses, and an empirical
account of when a calibrated threshold transfers.

Calibration of predictive scores is a separate line of work. Modern networks
are often miscalibrated \citep{guo2017calibration}, and a ranking score is not
a calibrated estimate of the quantity it orders. Context selection inherits
this problem: the utility model is trained to rank candidates
\citep{hashimoto-etal-2024-take,liu2022what,rubin2022learning}, so its output
needs an explicit calibration step before it can support a decision about
whether to route at all.

Negative transfer is the phenomenon that motivates the harm budget. Transfer
can degrade performance when source and target disagree
\citep{pan2010survey,torrey2010transfer,wang2019characterizing}, and recent
work mitigates it in sequential recommendation \citep{ant2025negative}. In our
setting the harmful event is local: adding one context increases the loss of
one prediction. The router cannot observe that event at deployment time, which
is what makes an abstention rule necessary.

Context selection itself is a large literature: retrieval and demonstration
selection \citep{brown2020language,min2022rethinking}, diversity and coverage
objectives \citep{carbonell1998use,gupta-etal-2023-coverage}, outcome-aware
valuations of individual candidates \citep{hashimoto-etal-2024-take}, and
learned routing over auxiliary inputs
\citep{jacobs1991adaptive,shazeer2017outrageously,rosenbaum2018routing}. These
methods rank candidates. Our question begins where the ranking ends: given a
ranked shortlist, should the system act, and how much of it can it act on
while keeping harmful selections within a budget?
